\scanpage{6}
donnant naissance au diagramme
\diagram{\begin{tikzcd}[ampersand replacement=\&]
{\unsure{\underline g|\underline{\mathfrak g}}}^{rss}\arrow[d,"f^{rss}"'] \& \underline{\mathscr L}^{rss}\arrow[l,"{{q}^{rss}}"']\arrow[d,"{\widetilde f'{}^{rss}}"]\\
{\unsure{J|I|\mathfrak J}}^{rss}=\underline t^{rss}/W \& J'=\underline t^{rss}\arrow[l].
\end{tikzcd}}

Supposons que $W$ opère \deleted{\unsure{triv.|trivial.}}\uline{librement} sur $\underline t^{rss}$ ; \inserted{il revient ici au \unsure{m\textsuperscript{me}|même} qu'il opère \deleted{\unsure{fidèl.|fidèlement}} \unsure{libr\textsuperscript{t}|librement|fidèlement} \deleted{\unsure{sur les}} sur les fibres géom de $\underline t$,} ce qui est le cas au moins si les car. résiduelles sont toutes $\ne 2$. Alors le diagramme précédent est cartésien ; \deleted{\unsure{il a un sens}}

Je dis que le morphisme $\underline g^{\unsure{rs|rss}}\longrightarrow J^{rss}$ se prolonge en un morphisme $\underline g\longrightarrow J$ (évidemment de façon unique), de sorte qu'on aura le diagramme commutatif
\diagram{\begin{tikzcd}[ampersand replacement=\&,column sep=small]
\underline g\arrow[d,"f"'] \& \& \underline{\mathscr L}\arrow[ll,bend right=13,"q"']\arrow[dl,"{\widetilde f'}"]\\
J=\underline t^{rss}/W \& J'=\underline t^{\deleted{rss}}/\arrow[l] \&
\end{tikzcd}}
\margin{\uline{NB} L'anneau affine de $J$ sur $S$ est le sous-anneau des invariants de l'anneau affine de $\underline g$ sur $S$, et $f$ décrit \unsure{le|la} \unsure{pb|prop.} universel d'envoyer $\underline g$ de façon invariante dans un schéma relatif \uline{affine} (restriction \unsure{possiblement|probablement} inutile, du moins si on prend des schémas relatifs séparés)}

Dém. de l'existence de $f$. On peut se ramener au cas $S=\operatorname{Spec}\mathbb Z$, $G$ déployé. \deleted{Alors il suffit de prouver que $f^{rss},q^{rss}$ se prolongent, et} Comme $\underline g$ est normal, $J$ affine, il suffit de le voir sur la \unsure{fibres|fibre} générique. Là c'est connu (Kostant ....)

Introduisons
\[
\underline g'=\underline g\times_J J'
\]
on trouve donc un diagramme
\diagram{\begin{tikzcd}[ampersand replacement=\&,column sep=large,execute at end picture={\draw[thin,gray] (gprime) ellipse (.40cm and .53cm);\draw[thin,gray] (gprime.south) to[bend right=12] (ell.south east) to[bend right=67] (ell.north) to[bend right=9] (gprime.north);}]
\underline g\arrow[d,"f"'] \& |[alias=gprime]|\underline g'\arrow[l]\arrow[d,"f'"'] \& |[alias=ell]|\underline{\mathscr L}\arrow[l,"p"]\arrow[ll,bend right=16,"q"']\arrow[dl,"{\bar f'}"]\arrow[r] \& D\\
J \& J'\arrow[l] \& \&
\end{tikzcd}}
avec $p$ morphisme projectif, $\bar f'$ morphisme lisse ; dans le cas favorable où $W$ opère \deleted{\unsure{fidèl\textsuperscript{t}|fidèlement}}\inserted{librement} sur $J'^{rss}$, i.e. fidèlement sur les fibres géom. de $\underline t$, alors $p$ \deleted{induit} est un isom birationnel, plus précisément il induit un isom $\mathscr L^{rss}\simeq\underline g'^{rss}$, où $\underline g^{rss}=\underline g'|\underline g^{rss}$.
% archive sheet number: 5
